\documentclass[11pt]{article}

\usepackage[a4paper,margin=2.4cm]{geometry}
\usepackage{amsmath,amssymb,amsthm,mathtools}
\usepackage{array}
\usepackage{booktabs}
\usepackage{longtable}
\usepackage{enumitem}
\usepackage[colorlinks=true,linkcolor=blue,citecolor=blue,urlcolor=blue]{hyperref}
\usepackage{xcolor}
\usepackage[export]{adjustbox}
\usepackage{microtype}
\usepackage{seqsplit}
\emergencystretch=2.5em

\title{The Inspector Column: Editable Code, Typing, and Outline\\
\large Implementation of Block-MiniJava's Three Program-View Tabs}
\author{Research note --- Block-MiniJava}
\date{2026-07-20}

\begin{document}
\maketitle

\begin{abstract}
The inspector column (\texttt{<aside id="code-column">}, \texttt{src/index.html:210--238})
is Block-MiniJava's dedicated "view the program as something other than
blocks" panel. It hosts three tabs that all read the \emph{same} live
workspace and never own a separate copy of the program: \textbf{Editable
Code} (a bidirectional MiniJava text $\leftrightarrow$ blocks editor),
\textbf{Typing} (the on-demand natural-deduction derivation viewer
documented separately in \texttt{typing-and-evaluation-rules.tex}), and
\textbf{Outline} (a collapsible tree of the block program's structure).
This note explains how each tab is implemented and kept synchronized with
workspace edits, with every mechanism traced to its file and line.
\end{abstract}

\tableofcontents

% =============================================================================
\section{The inspector column and its tab switch}

\begin{center}
\begin{longtable}{@{}p{3.6cm}p{11.6cm}@{}}
\toprule
\textbf{Element} & \textbf{Role} \\
\midrule
\endhead
\texttt{.inspector-tabs} &
  \texttt{index.html:219--223} --- three \texttt{<button role="tab">}s:
  \texttt{tab-code}, \texttt{tab-typing}, \texttt{tab-outline}, each with
  \texttt{data-panel} naming the panel it controls. \\
\texttt{.inspector-panel} &
  \texttt{index.html:225--236} --- three \texttt{<section role="tabpanel">}s:
  \texttt{panel-code}, \texttt{panel-typing}, \texttt{panel-outline}; the
  active one carries class \texttt{is-active}. \\
\texttt{InspectorPanel} type &
  \texttt{app.ts:70} --- \texttt{'code' | 'typing' | 'outline'}, the closed
  set \texttt{selectInspectorPanel} dispatches on. \\
\texttt{\seqsplit{selectInspectorPanel(panel)}} &
  \texttt{app.ts:353--366} --- the one function every tab click routes
  through: toggles \texttt{is-active}/\texttt{aria-selected} on the matching
  tab and panel pair, then triggers that panel's own on-activate render
  (\texttt{scheduleOutlineRender} for \texttt{'outline'},
  \texttt{scheduleTypingRender} for \texttt{'typing'} --- Editable Code
  needs no such call, since it stays synchronized continuously,
  \S\ref{sec:code}), and finally calls \texttt{requestLayoutResize(false)}
  since a newly-shown panel can change scrollbar/wrap width. \\
\bottomrule
\end{longtable}
\end{center}

Only \emph{Typing} and \emph{Outline} defer their rendering to tab
activation (\texttt{scheduleTypingRender}/\texttt{scheduleOutlineRender}
both \texttt{return} immediately if their panel is not the active one ---
see \texttt{typingPanel.ts:163--166} and \S\ref{sec:outline-render} ---
this is a cheap no-op check against needlessly rebuilding an invisible
DOM tree). \textbf{Editable Code} keeps no such guard: its textarea is the
one surface a user can type into at any time, so it must always reflect
(or be reflected onto) the workspace regardless of which tab is currently
showing.

% =============================================================================
\section{Editable Code}
\label{sec:code}

\texttt{src/core/ui/codeEditor.ts} (307 lines) implements a full
bidirectional editor: blocks $\to$ text is a pure, always-available
function (\texttt{generateMiniJava}); text $\to$ blocks is a debounced,
validated import that only ever fires from a real keystroke, never from a
workspace-originated change, so the two directions cannot feed each other
in a loop.

\subsection{Blocks $\to$ text}

\begin{center}
\begin{longtable}{@{}p{4.0cm}p{10.2cm}@{}}
\toprule
\textbf{Function} & \textbf{Role} \\
\midrule
\endhead
\texttt{\seqsplit{generateMiniJava(workspace)}} &
  \texttt{core/generator/minijavaGenerator.ts} (271 lines) --- walks the
  block tree from the \texttt{mj\_goal} block and emits MiniJava source
  text; the single source of truth both the read-only highlighted view and
  the editable textarea are kept in sync with. \\
\texttt{\seqsplit{highlightMiniJava(code)}} &
  \texttt{core/generator/highlighter.ts} (193 lines) --- a small MiniJava
  syntax highlighter (keywords, literals, comments) that renders the
  \emph{overlay} layer behind the transparent textarea (\S\ref{sec:code-dom}). \\
\texttt{\seqsplit{syncEditorFromWorkspace(workspace, code, message)}} &
  \texttt{codeEditor.ts:94--112} --- called by \texttt{app.ts}'s
  \texttt{updateCode()} (\texttt{app.ts:212--219}) after every workspace
  mutation. Two guards make it safe to call unconditionally: it is a no-op
  while the user is actively focused in the textarea
  (\texttt{document.activeElement === editor}, l.104), and it never
  overwrites text that already \emph{means} the same program
  (\S\ref{sec:structural-guard}). \\
\bottomrule
\end{longtable}
\end{center}

\subsection{The overlay DOM: line numbers, syntax color, and a transparent textarea}
\label{sec:code-dom}

\texttt{createEditorDom()} (\texttt{codeEditor.ts:218--258}) builds the
editor lazily, the first time \texttt{installEditableMiniJavaCodeEditor}
runs, by rearranging the read-only \texttt{<pre id="generated-code">}
element (\texttt{index.html:226}) that a headless/no-JS fallback would
otherwise just display statically:

\begin{enumerate}
  \item A new \texttt{.code-editor-pane} wraps a \texttt{.code-editor-line-numbers}
        \texttt{<pre>} (l.232--235, rebuilt by \texttt{renderLineNumbers},
        l.56--60) and the original \texttt{<pre class="code-view">} —
        the latter is repurposed as the syntax-\emph{highlight} layer, not
        removed.
  \item A real \texttt{<textarea id="generated-code-editor">} is stacked
        \emph{on top} of the highlight \texttt{<pre>} (l.238--246): the
        textarea is what the user actually types into and is visually
        transparent, while the \texttt{<pre>} underneath (kept in sync by
        \texttt{renderHighlight}, l.70--79) supplies the colored text the
        user sees through it.
  \item \texttt{syncHighlightScroll}/\texttt{syncHighlightMetrics}
        (l.50--54, l.64--68) keep the highlight layer's scroll offset and
        right-edge padding pixel-identical to the textarea's, including
        compensating for whichever browser's scrollbar happens to be
        present — a \texttt{ResizeObserver} (l.271--276) re-syncs on any
        panel resize.
\end{enumerate}
This "transparent textarea over a colored \texttt{<pre>}" construction is
the standard technique for a lightweight syntax-highlighting editor without
a full editor component (no CodeMirror/Monaco dependency): the browser's
native textarea supplies cursor, selection, undo, and IME handling for
free, and only the \emph{visual} rendering is swapped out underneath it.

\subsection{The import pipeline: debounce, parse, validate}

\begin{center}
\begin{longtable}{@{}p{4.0cm}p{10.2cm}@{}}
\toprule
\textbf{Function} & \textbf{Role} \\
\midrule
\endhead
\texttt{\seqsplit{scheduleEditorImport(workspace, options)}} &
  \texttt{codeEditor.ts:191--200} --- fired on every \texttt{input} event
  (wired at l.266). Re-renders the highlight layer and status text
  immediately (so the view never looks stale), then debounces the actual
  import by \texttt{TEXT\_IMPORT\_DELAY\_MS = 650}ms (l.6) via
  \texttt{window.setTimeout}, restarting the timer on every keystroke ---
  so a workspace rebuild only happens once typing pauses, not on every
  character. \\
\texttt{\seqsplit{applyEditorTextToWorkspace(...)}} &
  \texttt{codeEditor.ts:157--189} --- parses the trimmed textarea contents
  with \texttt{parseMiniJavaTextToWorkspaceState}
  (\texttt{core/parser/minijavaTextParser.ts}, 853 lines); on a parse
  error, catches \texttt{MiniJavaTextParseError} and reports it in the
  status line and via \texttt{options.onImportError} rather than throwing
  through to the caller — an in-progress, syntactically incomplete edit is
  an expected state, not a bug. \\
\texttt{\seqsplit{importEditorTextToWorkspace(...)}} &
  \texttt{codeEditor.ts:138--155} --- the actual rebuild: sets a
  re-entrancy guard (\S\ref{sec:suppression}), \texttt{workspace.clear()},
  loads the parsed \texttt{Blockly.serialization.workspaces} state, and
  calls \texttt{workspace.cleanUp()} to re-lay-out top blocks. \\
\bottomrule
\end{longtable}
\end{center}

\subsection{The structural-equality guard: never delete what the user just typed}
\label{sec:structural-guard}

Two problems have to be solved together for a bidirectional editor to feel
safe to type in: (a) the generator's output is not necessarily
byte-identical to what the user wrote (whitespace, comments, formatting
choices), and (b) a full \texttt{workspace.clear()} + rebuild on every
keystroke would destroy block ids, undo history, and manual layout for a
change that produces the \emph{same program}. \texttt{codeEditor.ts} solves
both with one idea: compare \emph{parsed structure}, not text.

\begin{center}
\begin{longtable}{@{}p{4.4cm}p{9.8cm}@{}}
\toprule
\textbf{Function} & \textbf{Role} \\
\midrule
\endhead
\texttt{\seqsplit{workspaceStructureSignature(ws)}} &
  \texttt{codeEditor.ts:116--122} --- re-generates the workspace's own code
  and re-parses \emph{that}, then \texttt{JSON.stringify}s the resulting
  serialization state: the workspace's canonical structural signature. \\
\texttt{\seqsplit{editorTextMatchesCode(text, code)}} &
  \texttt{codeEditor.ts:128--136} --- true iff the editor's raw text parses
  to the identical structure as \texttt{code} (an exact string match short-circuits
  first, l.129). Used by \texttt{syncEditorFromWorkspace}
  (\texttt{blocks$\to$text} direction, l.107) to decide whether to
  overwrite the textarea at all: comment/whitespace-only edits are left
  alone. \\
\texttt{\seqsplit{applyEditorTextToWorkspace}}'s skip check &
  \texttt{codeEditor.ts:176--179} (\texttt{text$\to$blocks} direction) ---
  compares the freshly parsed state against
  \texttt{workspaceStructureSignature(workspace)} \emph{before} rebuilding;
  on a match it reports "already matches the blocks" and returns without
  touching the workspace at all. \\
\bottomrule
\end{longtable}
\end{center}

\subsection{Breaking the sync feedback loop}
\label{sec:suppression}

A naive bidirectional sync is prone to an infinite (or at least
user-visible flicker) loop: importing text mutates the workspace, the
workspace's change listener calls \texttt{updateCode()}
(\texttt{app.ts:1122}), which calls \texttt{syncEditorFromWorkspace}, which
could re-touch the very textarea that triggered the import in the first
place. \texttt{suppressWorkspaceSyncUntil} (\texttt{codeEditor.ts:35}) is a
simple timestamp guard: \texttt{importEditorTextToWorkspace} sets it 1500ms
into the future (l.146) before mutating the workspace, and
\texttt{syncEditorFromWorkspace} bails out immediately (l.100) while
"now" is still before that deadline. Combined with the
\texttt{document.activeElement === editor} check (\S\ref{sec:code}),
this means: while the user is typing, nothing else may touch the textarea;
immediately after an import, nothing may touch it either, until the
workspace's own settling mutations (autosave, required-block enforcement)
have had time to finish.

\subsection{Other editor entry points}

\texttt{EditableMiniJavaCodeEditor}'s public surface
(\texttt{codeEditor.ts:21--27}) exposes \texttt{loadText} in addition to
the input-driven path: called from \textbf{Open} (\texttt{app.ts:766}) to
force-import a loaded \texttt{.java}/\texttt{.txt} file even when its
structure happens to already match the current blocks (the file's
\emph{label} must still update, l.291--302, hence
\texttt{allowSkipWhenUnchanged = false} there). \texttt{resize()} is called
by the shared layout coordinator (\texttt{app.ts}'s
\texttt{requestLayoutResize}, cross-referenced in
\texttt{ui-design-and-layout.tex}) whenever the inspector column changes
width.

% =============================================================================
\section{Typing}

The Typing tab's implementation --- \texttt{deriveProgram} in
\texttt{core/types/typeChecker.ts} and the renderer in
\texttt{core/ui/typingPanel.ts} --- is documented in full in
\texttt{typing-and-evaluation-rules.tex} (the typing-rule inference-rule
reconstruction) and was the subject of the project's most recent commit.
This note gives only the piece that document does not: how it fits into
the shared inspector-tab machinery of \S1.

\begin{center}
\begin{longtable}{@{}p{4.2cm}p{10.0cm}@{}}
\toprule
\textbf{Symbol} & \textbf{Role} \\
\midrule
\endhead
\texttt{\seqsplit{scheduleTypingRender()}} &
  \texttt{typingPanel.ts:163--171} --- rAF-debounced; \texttt{returns}
  immediately unless \texttt{\#panel-typing} currently carries
  \texttt{is-active} (l.165), so a background workspace edit while another
  tab is open costs nothing. Called both from
  \texttt{selectInspectorPanel('typing')} (\S1) and from the workspace's
  own change listener (\texttt{app.ts:1109}) on every non-viewport
  mutation. \\
\texttt{\seqsplit{initTypingPanel(getWorkspace)}} &
  \texttt{typingPanel.ts:173--180} --- wires the \texttt{\#typing-method-select}
  dropdown's \texttt{change} listener and stores the workspace getter;
  called once from \texttt{app.ts:1416}. \\
\texttt{\seqsplit{deriveProgram(workspace)}} &
  \texttt{core/types/typeChecker.ts:495--497} --- the same walk that
  produces the Problems tab's diagnostics (see \texttt{bottom-panel-sync.tex}),
  so the two views can never disagree. \\
\bottomrule
\end{longtable}
\end{center}

% =============================================================================
\section{Outline}
\label{sec:outline}

The Outline tab renders the block program as a plain indented tree ---
every top-level block, recursively, down to its structural children ---
each row click-to-locate to the corresponding block in the workspace. It
is implemented entirely inside \texttt{app.ts} (no separate module).

\subsection{Rendering}
\label{sec:outline-render}

\begin{center}
\begin{longtable}{@{}p{4.2cm}p{10.0cm}@{}}
\toprule
\textbf{Function} & \textbf{Role} \\
\midrule
\endhead
\texttt{\seqsplit{scheduleOutlineRender()}} &
  \texttt{app.ts:368--374} --- rAF-debounced (mirrors
  \texttt{scheduleTypingRender}'s pattern, though this one does not itself
  check tab visibility --- the caller at \texttt{selectInspectorPanel}
  gates on the panel becoming active, and the workspace listener
  (\texttt{app.ts:1108}) fires it unconditionally on every edit regardless
  of which tab is showing, since a full outline rebuild is cheap enough
  not to need the same guard). \\
\texttt{renderOutline()} &
  \texttt{app.ts:376--427} --- clears \texttt{\#program-outline}
  (\texttt{index.html:235}); if the workspace has no top blocks, shows an
  \texttt{.outline-empty} placeholder (l.382--388); otherwise walks every
  top block via the local \texttt{addBlock(block, depth)} closure
  (l.390--424). \\
\texttt{\seqsplit{addBlock(block, depth)}} &
  \texttt{app.ts:390--424} (inline closure) --- recursive tree walk using
  Blockly's own \texttt{block.getChildren(true)} (structural children,
  l.391) as the recursion source --- \textbf{no separate outline data
  structure is built or maintained}; the block tree \emph{is} the outline,
  walked fresh on every render. Each row is a \texttt{button.outline-item}
  with a CSS custom property \texttt{--outline-depth} driving indentation
  (l.396), a disclosure glyph (a caret if the block has children, a
  centered dot otherwise, l.400--403), the block's own \texttt{toString(54,
  '…')} rendering truncated to 54 characters as the row label
  (l.407--408), and the block's type with the \texttt{mj\_} prefix and
  underscores stripped as a small trailing tag (l.410--412). \\
\bottomrule
\end{longtable}
\end{center}

\subsection{Click-to-locate}

Each row's click handler (\texttt{app.ts:415--421}) is the same
three-call idiom used everywhere else in the IDE that links a rendered
artifact back to its source block (Problems, CESK, A~vs~B --- see the
companion \texttt{bottom-panel-sync.tex}):
\texttt{workspace.getBlockById(item.dataset.blockId)}, then
\texttt{workspace.centerOnBlock(id)}, then
\texttt{Blockly.common.setSelected(block)} (which, unlike
\texttt{BlockSvg.select()}, also clears whatever was previously selected).
The Outline row itself carries no separate "locate" affordance --- the
\emph{whole row} is the button, since unlike Problems (which also needs
room for a multi-line diagnostic message) an outline entry has nothing
else to click.

\subsection{Why Outline needs no staleness guard}

Unlike the CESK/Compare/Rewrite steppers (which snapshot a
\texttt{MachineState} or a copied block tree that can drift out of sync
with a live-edited workspace, requiring an explicit \texttt{stale} flag —
see \texttt{stepperPanel.ts:46} and the companion CESK note), the Outline
and Typing tabs hold \textbf{no persistent state of their own} between
renders: each call to \texttt{renderOutline}/\texttt{deriveProgram} reads
the live workspace fresh. There is nothing to go stale, and hence no
"Program changed --- press Load to restart" banner in either tab, unlike
the machine steppers.

% =============================================================================
\section{Cross-cutting: the provenance idiom}

All three tabs, and every other program-view surface in the IDE, converge
on the same two- or three-line idiom to link a rendered UI element back to
the block that produced it:

\begin{center}
\begin{longtable}{@{}p{3.4cm}p{10.8cm}@{}}
\toprule
\textbf{Tab} & \textbf{Locate function} \\
\midrule
\endhead
Editable Code &
  N/A --- the editor \emph{is} a textual view of the whole program at
  once; there is no per-line click-to-locate, only the debounced
  text$\to$blocks import of \S\ref{sec:code}. \\
Typing &
  Every judgement button calls \texttt{locate(blockId)}
  (\texttt{typingPanel.ts:26--30}), which forwards to
  \texttt{locateBlock} from \texttt{typeDiagnostics.ts} (shared with the
  Problems tab, \texttt{typeDiagnostics.ts:37--44}). \\
Outline &
  Inline in \texttt{addBlock}'s click handler
  (\texttt{app.ts:415--421}), using \texttt{centerOnBlock} +
  \texttt{Blockly.common.setSelected} directly. \\
\bottomrule
\end{longtable}
\end{center}

The full inventory of this pattern across the \emph{bottom} panel
(Problems, CESK, A~vs~B, and the two tabs that deliberately \emph{don't}
sync back to the main workspace) is catalogued in the companion note
\texttt{bottom-panel-sync.tex}.

% =============================================================================
\section*{Note on scope}
\addcontentsline{toc}{section}{Note on scope}
This note covers the three \emph{inspector column} tabs. The seven
\emph{bottom panel} (viz-dock) tabs --- Problems, Output, Call-by-Structure,
Call-by-Value, CESK, A~vs~B, Rewrite --- are covered in
\texttt{bottom-panel-sync.tex}. The Typing tab's typing-rule content is
covered in \texttt{typing-and-evaluation-rules.tex}. The overall workbench
shell (activity bar, perspectives, panel visibility/resizing) that the
inspector column sits inside is covered in \texttt{ui-design-and-layout.tex}.

\end{document}
